\chapter*{Résumé}
\markboth{Resume}{}
\addstarredchapter{Resume}

Ce mémoire présente un projet réalisé dans le cadre du module DNF2ED12 (Master 1, mention Informatique, parcours Big Data et Fouille de Données, Université Paris 8) : une application qui extrait automatiquement les informations contenues dans des documents d'identité (carte nationale d'identité, passeport, certificat de scolarité, titre de séjour) à partir d'une simple photo, puis vérifie que ces informations correspondent bien à celles d'une base de référence.

Le système repose sur trois étapes : un prétraitement de l'image (redressement, amélioration du contraste, débruitage), une reconnaissance optique de caractères avec le moteur Tesseract, puis une comparaison floue avec les données de référence (fichier Excel) à l'aide de la bibliothèque rapidfuzz. Une interface web développée avec Streamlit permet à un utilisateur de charger un document et de visualiser les résultats.

Une fois l'application fonctionnelle, un travail spécifique a été mené pour améliorer la précision de la vérification, jugée insuffisante en l'état. Un banc d'essai (benchmark) a été construit à partir d'images synthétiques dégradées de sept manières différentes (rotation, flou, bruit, basse résolution, faible contraste...), avec une vérité terrain connue pour chaque champ. Ce banc d'essai a permis de mesurer objectivement les progrès, et surtout de révéler trois défauts cachés du pipeline : l'absence de limite de temps sur l'OCR (qui pouvait bloquer indéfiniment sur une image très bruitée), une erreur de signe dans la correction d'inclinaison des images (qui doublait l'inclinaison au lieu de la corriger), et une mauvaise hypothèse sur la convention d'angle renvoyée par une fonction d'OpenCV. La correction de ces trois défauts a fait passer la précision globale mesurée de 57,1\,\% à 79,1\,\% des champs correctement vérifiés, soit un gain de 22 points. Une suite de 41 tests unitaires garantit la non-régression de ces correctifs. Le système a ensuite été étendu à un quatrième type de document (titre de séjour), ce qui a permis d'identifier et de corriger cinq défauts supplémentaires liés à l'analyse du texte OCR (ordre de lecture, fond de sécurité imprimé, collision entre champs, format de carte récent, lecture de zone MRZ), validés sur des documents réels mais non encore couverts par le banc d'essai chiffré ni par les tests automatisés.

\textbf{Mots-clés :} OCR, Tesseract, traitement d'image, vérification de documents d'identité, comparaison floue, benchmark, tests unitaires.
